% ============================================================
%  Chapitre 00 — Analyse dimensionnelle et remise à niveau
%  BTS FED option GCF — Cours complet
%  Créé le 2026-04-07
% ============================================================

\chapter{Analyse dimensionnelle et remise à niveau mathématique}

% ============================================================
\section{Objectifs pédagogiques}
% ============================================================

\subsection*{Compétences GCF mobilisées}

\begin{itemize}
  \item \textbf{C1-1} — Exploiter des sources d'information : récupérer et interpréter des
        données techniques (fiches produit, tableaux, courbes).
  \item \textbf{C3-1} — Analyser un système : identifier les paramètres et les grandeurs
        physiques pertinentes.
  \item \textbf{C3-3} — Dimensionner : maîtriser les transformations d'équations et les
        conversions d'unités.
  \item \textbf{C7-4} — Conduire des mesures : interpréter les résultats en cohérence avec
        les unités employées.
\end{itemize}

\subsection*{Savoirs associés GCF}

\begin{itemize}
  \item \textbf{S1} — Mathématiques appliquées : calcul littéral, résolution d'équations,
        exploitation de graphiques. (Niveau GCF\,: 3)
  \item \textbf{S4} — Sciences physiques : grandeurs, unités SI et dérivées, analyse
        dimensionnelle. (Niveau GCF\,: 2)
\end{itemize}

% ============================================================
\section{Introduction}
% ============================================================

L'analyse dimensionnelle est une méthode fondamentale en sciences de l'ingénieur permettant
de vérifier la cohérence physicienne d'une équation ou d'un calcul. En génie climatique, on
rencontre fréquemment des formules complexes mettant en jeu des puissances (W), des débits
(m³/h, kg/s), des pressions (Pa, bar), des températures (K, °C)…

Ce chapitre fournit les outils pour :
\begin{itemize}
  \item \textbf{Vérifier l'homogénéité} d'une formule avant de l'appliquer ;
  \item \textbf{Convertir des unités} avec assurance et rigueur ;
  \item \textbf{Transformer des équations} pour exprimer une inconnue en fonction d'autres
        paramètres ;
  \item \textbf{Exploiter des graphiques et abaques} : lecture de courbes, interpolation,
        extrapolation dans un réseau d'isobares ou d'isothermes.
\end{itemize}

La maîtrise de ces concepts garantit la fiabilité des dimensionnements et la détection rapide
d'erreurs.

% ============================================================
\section{Système international d'unités (SI)}
% ============================================================

\subsection{Les sept unités de base}

Le Système International (SI) repose sur sept unités de base, dont sont dérivées toutes les
autres :

\begin{table}[h]
  \centering
  \sffamily
  \caption{Unités de base du SI}
  \label{tab:unite-base}
  \begin{tabular}{|l|c|l|}
    \hline
    \textbf{Grandeur} & \textbf{Symbole} & \textbf{Unité SI} \\
    \hline
    Longueur & $\ell$, $L$, $d$ & \si{\meter} (m) \\
    Masse & $m$ & \si{\kilogram} (kg) \\
    Temps & $t$ & \si{\second} (s) \\
    Intensité de courant & $I$ & \si{\ampere} (A) \\
    Température thermodynamique & $T$ & \si{\kelvin} (K) \\
    Quantité de matière & $n$ & \si{\mole} (mol) \\
    Intensité lumineuse & $I_v$ & \si{\candela} (cd) \\
    \hline
  \end{tabular}
\end{table}

En génie climatique, les trois premières unités (longueur, masse, temps) sont prédominantes,
associées à la température thermodynamique.

\subsection{Unités dérivées courantes en CVC}

De nombreuses grandeurs dérivées des sept unités de base sont d'usage courant. Le tableau
ci-dessous en liste les plus importantes :

\begin{table}[h]
  \centering
  \sffamily
  \caption{Unités dérivées du SI usuelles en génie climatique}
  \label{tab:unite-derivee}
  \begin{tabular}{|l|c|c|l|}
    \hline
    \textbf{Grandeur} & \textbf{Symbole} & \textbf{Unité SI} & \textbf{Expression en base} \\
    \hline
    Vitesse & $v$ & \si{\meter\per\second} (m/s) & $\mathrm{m \cdot s^{-1}}$ \\
    Accélération & $a$ & \si{\meter\per\second\squared} (m/s²) & $\mathrm{m \cdot s^{-2}}$ \\
    Force & $F$ & \si{\newton} (N) & $\mathrm{kg \cdot m \cdot s^{-2}}$ \\
    Pression & $P$ & \si{\pascal} (Pa) & $\mathrm{kg \cdot m^{-1} \cdot s^{-2}}$ \\
    Énergie, travail & $E$, $W$ & \si{\joule} (J) & $\mathrm{kg \cdot m^2 \cdot s^{-2}}$ \\
    Puissance & $\dot{Q}$, $P$ & \si{\watt} (W) & $\mathrm{kg \cdot m^2 \cdot s^{-3}}$ \\
    Débit massique & $\dot{m}$ & \si{\kilogram\per\second} (kg/s) & $\mathrm{kg \cdot s^{-1}}$ \\
    Débit volumique & $\dot{V}$ & \si{\meter\cubed\per\second} (m³/s) & $\mathrm{m^3 \cdot s^{-1}}$ \\
    \hline
  \end{tabular}
\end{table}

% ============================================================
\section{Analyse dimensionnelle — Principe d'homogénéité}
% ============================================================

\subsection{Définition}

\begin{definition}
  Une équation physique est \textbf{dimensionnellement homogène} si, et seulement si, tous
  les termes de l'équation ont la même dimension. Chaque terme doit s'exprimer dans les mêmes
  unités (ou équivalentes via les unités dérivées).
\end{definition}

\textbf{Principe cardinal} : \emph{On ne peut ajouter ou soustraire que des grandeurs de même
dimension. On ne peut comparer que des grandeurs de même dimension.}

\subsection{Méthode — Vérifier l'homogénéité}

Soit l'équation de bilan thermique d'un échangeur :
\begin{equation}
  \dot{Q} = \dot{m}_f \cdot c_p \cdot \Delta T
  \label{eq:bilan-thermal}
\end{equation}

où :
\begin{itemize}
  \item $\dot{Q}$ : puissance thermique échangée (\si{\watt}, W)
  \item $\dot{m}_f$ : débit massique du fluide (\si{\kilogram\per\second}, kg/s)
  \item $c_p$ : capacité thermique massique (\si{\joule\per\kilogram\per\kelvin}, J/(kg·K))
  \item $\Delta T$ : différence de température (\si{\kelvin}, K)
\end{itemize}

\vspace{0.5cm}

\textbf{Analyse de dimension} :

\begin{align*}
  [\dot{Q}] &= \mathrm{W} = \mathrm{kg \cdot m^2 \cdot s^{-3}} \\
  [\dot{m}_f \cdot c_p \cdot \Delta T] &= [\mathrm{kg/s}] \times [\mathrm{J/(kg \cdot K)}] \times [
  \mathrm{K}] \\
  &= \mathrm{(kg \cdot s^{-1})} \times \mathrm{(kg^{-1} \cdot m^2 \cdot s^{-2} \cdot K^{-1})} \times
  \mathrm{K} \\
  &= \mathrm{kg \cdot s^{-1} \cdot kg^{-1} \cdot m^2 \cdot s^{-2} \cdot K \cdot K^{-1}} \\
  &= \mathrm{m^2 \cdot s^{-3}} = \mathrm{W}
\end{align*}

$\Rightarrow$ L'équation \eqref{eq:bilan-thermal} est dimensionnellement homogène. ✓

\subsection*{Exemple 1 — Détection d'erreur}

Supposons une formule proposée (erronée) :
\begin{equation*}
  P_{\text{pompe}} = \dot{m} + \Delta P
\end{equation*}

où $P_{\text{pompe}}$ est en W, $\dot{m}$ en kg/s et $\Delta P$ en Pa.

\[
  [\mathrm{W}] \stackrel{?}{=} [\mathrm{kg/s}] + [\mathrm{Pa}]
\]

Les dimensions ne correspondent pas :
\begin{itemize}
  \item Membre de gauche : $\mathrm{kg \cdot m^2 \cdot s^{-3}}$
  \item Membre de droite (1) : $\mathrm{kg \cdot s^{-1}}$
  \item Membre de droite (2) : $\mathrm{kg \cdot m^{-1} \cdot s^{-2}}$
\end{itemize}

$\Rightarrow$ La formule est \textbf{fausse}. La bonne formule est : $P_{\text{pompe}} = \dot{V}
\times \Delta P$

% ============================================================
\section{Conversions d'unités}
% ============================================================

\subsection{Principes généraux}

Une conversion d'unité consiste à exprimer une grandeur dans une unité différente, sans en
modifier la valeur physique. Trois approches sont possibles :

\begin{enumerate}
  \item \textbf{Utiliser un facteur de conversion} (rapport dimensionnel)
  \item \textbf{Passer par l'unité SI} de référence
  \item \textbf{Utiliser des tableaux prédéfinis}
\end{enumerate}

\subsection{Conversions courantes en génie climatique}

\subsubsection{Puissance et énergie}

\begin{table}[h]
  \centering
  \sffamily\small
  \caption{Conversions courantes — Énergie et puissance}
  \label{tab:conv-energie}
  \begin{tabular}{|l|c|l|}
    \hline
    \textbf{De} & \textbf{Vers} & \textbf{Facteur de conversion} \\
    \hline
    \si{\watt} (W) & \si{\kilowatt} (kW) & $\times 10^{-3}$ \\
    \si{\joule} (J) & \si{\kilo\watt\hour} (kWh) & $\times 2{,}778 \times 10^{-7}$ \\
    kWh & \si{\joule} (J) & $\times 3{,}6 \times 10^6$ \\
    Wh & kWh & $\times 10^{-3}$ \\
    cal & \si{\joule} (J) & $\times 4{,}1868$ \\
    kcal & \si{\joule} (J) & $\times 4186{,}8$ \\
    \hline
  \end{tabular}
\end{table}

\textbf{Exemple — Conversion W $\leftrightarrow$ kW} :

Une chaudière gaz produit une puissance de $35~\mathrm{kW}$.

\[
  P = 35 \text{ kW} = 35 \times 10^3 \text{ W} = 35\,000 \text{ W}
\]

\subsubsection{Pression}

\begin{table}[h]
  \centering
  \sffamily\small
  \caption{Conversions courantes — Pression}
  \label{tab:conv-pression}
  \begin{tabular}{|l|c|l|}
    \hline
    \textbf{De} & \textbf{Vers} & \textbf{Facteur de conversion} \\
    \hline
    \si{\pascal} (Pa) & bar & $\times 10^{-5}$ \\
    bar & \si{\pascal} (Pa) & $\times 10^5$ \\
    bar & kgf/m² & $\times 10\,197{,}16$ \\
    \si{\pascal} (Pa) & mH\textsubscript{2}O & $\times 1{,}02 \times 10^{-4}$ \\
    mH\textsubscript{2}O & \si{\pascal} (Pa) & $\times 9\,810$ \\
    bar & mH\textsubscript{2}O & $\times 10{,}2$ \\
    \hline
  \end{tabular}
\end{table}

\textbf{Exemple — Conversion Pa $\leftrightarrow$ bar} :

Un détendeur hydrodynamique supporte une perte de charge $\Delta P = 25\,000~\mathrm{Pa}$.

\[
  \Delta P = 25\,000 \text{ Pa} = 25\,000 \times 10^{-5} \text{ bar} = 0{,}25 \text{ bar}
\]

\subsubsection{Débit volumique et débit massique}

\begin{table}[h]
  \centering
  \sffamily\small
  \caption{Conversions courantes — Débits}
  \label{tab:conv-debit}
  \begin{tabular}{|l|c|l|}
    \hline
    \textbf{De} & \textbf{Vers} & \textbf{Facteur de conversion} \\
    \hline
    \si{\meter\cubed\per\hour} (m³/h) & \si{\meter\cubed\per\second} (m³/s) & $\times 1/3\,600 \approx 2{,}778 \times 10^{-4}$ \\
    \si{\meter\cubed\per\second} (m³/s) & m³/h & $\times 3\,600$ \\
    L/s & m³/h & $\times 3{,}6$ \\
    m³/h & L/s & $\times 1/3{,}6 \approx 0{,}2778$ \\
    m³/h & L/min & $\times 1000/60 \approx 16{,}67$ \\
    \hline
  \end{tabular}
\end{table}

\textbf{Relation débit massique / débit volumique} :

\[
  \dot{m} = \rho \cdot \dot{V}
\]

Avec $\rho$ densité (kg/m³).

\textbf{Exemple — Conversion m³/h $\leftrightarrow$ m³/s} :

Un extracteur d'air a un débit volumique $\dot{V} = 800~\mathrm{m^3/h}$.

\begin{align*}
  \dot{V} &= 800 \text{ m³/h} = 800 \times \frac{1}{3\,600} \text{ m³/s} \\
  &= 800 \times 2{,}778 \times 10^{-4} \text{ m³/s} \\
  &= 0{,}222 \text{ m³/s}
\end{align*}

Avec une densité air $\rho = 1{,}2~\mathrm{kg/m^3}$, le débit massique est :

\[
  \dot{m} = 1{,}2 \times 0{,}222 = 0{,}267 \text{ kg/s}
\]

\subsubsection{Température}

\begin{table}[h]
  \centering
  \sffamily\small
  \caption{Conversions courantes — Température}
  \label{tab:conv-temp}
  \begin{tabular}{|l|c|l|}
    \hline
    \textbf{De} & \textbf{Vers} & \textbf{Formule} \\
    \hline
    Celsius (°C) & Kelvin (K) & $T\,[\mathrm{K}] = \theta\,[\mathrm{°C}] + 273{,}15$ \\
    Kelvin (K) & Celsius (°C) & $\theta\,[\mathrm{°C}] = T\,[\mathrm{K}] - 273{,}15$ \\
    Celsius (°C) & Fahrenheit (°F) & $T\,[\mathrm{°F}] = \theta\,[\mathrm{°C}] \times \frac{9}{5} + 32$ \\
    \hline
  \end{tabular}
\end{table}

\textbf{Remarque importante} : Il convient de distinguer \textbf{température absolue} (K, utilisée
en thermodynamique) et \textbf{écart de température} (K ou °C, utilisé pour les différences).

Une température $\theta = 20~\mathrm{°C}$ correspond à $T = 293{,}15~\mathrm{K}$.

Un écart $\Delta T = 5~\mathrm{K} \leftrightarrow \Delta \theta = 5~\mathrm{°C}$.

% ============================================================
\section{Transformation d'équations et résolution}
% ============================================================

\subsection{Principes de manipulation algébrique}

La résolution d'équations repose sur l'application rigoureuse d'opérations mathématiques
identiques aux deux membres de l'égalité.

\textbf{Règles fondamentales} :
\begin{itemize}
  \item Additionner/soustraire le même terme aux deux membres
  \item Multiplier/diviser par le même terme non nul aux deux membres
  \item Élever à la même puissance (avec attention aux racines)
  \item Appliquer une fonction (log, exp, sin, …) aux deux membres
\end{itemize}

\subsection{Exemple 1 — Équation linéaire simple}

\textbf{Problème} : Exprimer la température de sortie d'un radiateur en fonction des autres
paramètres.

Bilan énergétique du radiateur :

\[
  \dot{Q} = \dot{m} \cdot c_p \cdot (T_{\mathrm{in}} - T_{\mathrm{out}})
\]

Chercher $T_{\mathrm{out}}$ :

\begin{align*}
  \dot{Q} &= \dot{m} \cdot c_p \cdot (T_{\mathrm{in}} - T_{\mathrm{out}}) \\
  \frac{\dot{Q}}{\dot{m} \cdot c_p} &= T_{\mathrm{in}} - T_{\mathrm{out}} \quad \text{(diviser par
  $\dot{m} \cdot c_p$)} \\
  T_{\mathrm{out}} &= T_{\mathrm{in}} - \frac{\dot{Q}}{\dot{m} \cdot c_p}
\end{align*}

\subsection{Exemple 2 — Équation non linéaire}

\textbf{Problème} : Exprimer le diamètre d'une conduite en fonction du débit et de la vitesse.

Formule de continuité :

\[
  \dot{V} = v \cdot A = v \cdot \pi \cdot \frac{d^2}{4}
\]

Chercher $d$ :

\begin{align*}
  \dot{V} &= v \cdot \pi \cdot \frac{d^2}{4} \\
  \frac{4 \cdot \dot{V}}{v \cdot \pi} &= d^2 \quad \text{(isoler $d^2$)} \\
  d &= \sqrt{\frac{4 \cdot \dot{V}}{v \cdot \pi}} \\
  d &= 2 \cdot \sqrt{\frac{\dot{V}}{v \cdot \pi}}
\end{align*}

\subsection{Exemple 3 — Équation avec logarithme}

\textbf{Problème} : Exprimer le nombre d'unités de transfert en fonction de l'efficacité d'un
échangeur.

Relation DTLM (Différence de Température Logarithmique Moyenne) :

\[
  \dot{Q} = U \cdot A \cdot \Delta T_{\mathrm{lm}}
\]

avec $\Delta T_{\mathrm{lm}} = \dfrac{\Delta T_1 - \Delta T_2}{\ln(\Delta T_1 / \Delta T_2)}$

Pour certaines configurations, on utilise : $\varepsilon = 1 - \exp(-\mathrm{NUT})$

Chercher $\mathrm{NUT}$ en fonction de $\varepsilon$ :

\begin{align*}
  \varepsilon &= 1 - \exp(-\mathrm{NUT}) \\
  \varepsilon - 1 &= - \exp(-\mathrm{NUT}) \\
  1 - \varepsilon &= \exp(-\mathrm{NUT}) \quad \text{(retirer les signes)} \\
  \ln(1 - \varepsilon) &= -\mathrm{NUT} \quad \text{(appliquer ln)} \\
  \mathrm{NUT} &= -\ln(1 - \varepsilon)
\end{align*}

% ============================================================
\section{Exploitation de graphiques et abaques}
% ============================================================

\subsection{Lecture directe}

Certaines données sont fournies graphiquement (courbes, nomogrammes, tables). Savoir extraire
une valeur nécessite :

\begin{itemize}
  \item \textbf{Identifier les axes} : abscisse (x), ordonnée (y), unités
  \item \textbf{Localiser le point} correspondant à la donnée cherchée
  \item \textbf{Lire la valeur} avec la précision appropriée à l'échelle
  \item \textbf{Respecter les unités} du graphique
\end{itemize}

\textbf{Exemple — Propriétés de l'eau} :

Un abaque standard donne $\rho_{\mathrm{eau}}$ et $\nu_{\mathrm{eau}}$ en fonction de la
température.

À $T = 60°C$ : on lit $\rho = 983{,}2~\mathrm{kg/m^3}$ et $\nu = 4{,}74 \times
10^{-7}~\mathrm{m^2/s}$.

\subsection{Interpolation linéaire}

Lorsque le point cherché se situe \textit{entre} deux points connus, on fait une interpolation
linéaire :

\[
  y = y_1 + (y_2 - y_1) \times \frac{x - x_1}{x_2 - x_1}
\]

\textbf{Exemple — Déduction d'une valeur intermédiaire} :

Tableau de viscosité cinématique de l'air (à 1 atm) :

\begin{center}
  \sffamily
  \begin{tabular}{|c|c|}
    \hline
    $T\,[\mathrm{°C}]$ & $\nu\,[\times 10^{-5}~\mathrm{m^2/s}]$ \\
    \hline
    0 & 13{,}28 \\
    20 & 15{,}89 \\
    40 & 18{,}97 \\
    \hline
  \end{tabular}
\end{center}

Chercher $\nu$ à $T = 25°C$ :

\begin{align*}
  \nu(25) &= 15{,}89 + (18{,}97 - 15{,}89) \times \frac{25 - 20}{40 - 20} \\
  &= 15{,}89 + 3{,}08 \times \frac{5}{20} \\
  &= 15{,}89 + 0{,}77 \\
  &= 16{,}66 \times 10^{-5}~\mathrm{m^2/s}
\end{align*}

\subsection{Lecture sur un réseau de courbes (diagramme de Mollier, …)}

Les diagrammes complexes (psychrométrique, Mollier, Ashrae) permettent la détermination
simultanée de plusieurs propriétés.

\textbf{Méthode générale} :
\begin{enumerate}
  \item Localiser deux propriétés connues (ex : T et HR)
  \item Suivre les isoçorbes/isobares/isothermes appropriées
  \item Lire les propriétés secondaires au point d'intersection (ex : humidité absolue,
        enthalpie)
  \item Exprimer le résultat dans l'unité du diagramme
\end{enumerate}

% ============================================================
\section{Exercices résolus}
% ============================================================

\subsection*{Exercice 1 — Vérification d'homogénéité}

\textbf{Énoncé} : Vérifier l'homogénéité dimensionnelle de la formule de perte de charge en
conduite :

\[
  \Delta P = \lambda \cdot \frac{\ell}{d} \cdot \frac{\rho \cdot v^2}{2}
\]

où $\lambda$ est un coefficient sans dimension (coefficient de frottement), $\ell$ la longueur
(m), $d$ le diamètre (m), $\rho$ la masse volumique (kg/m³), $v$ la vitesse (m/s).

\textbf{Solution} :

\begin{align*}
  [\Delta P] &= \mathrm{Pa} = \mathrm{kg \cdot m^{-1} \cdot s^{-2}} \\
  [\lambda \cdot \frac{\ell}{d} \cdot \frac{\rho \cdot v^2}{2}] &= [1] \times [\mathrm{m/m}]
  \times \frac{[\mathrm{kg/m^3}] \times [\mathrm{m/s}]^2}{[1]} \\
  &= [\mathrm{kg \cdot m^{-3} \cdot m^2 \cdot s^{-2}}] \\
  &= [\mathrm{kg \cdot m^{-1} \cdot s^{-2}}] \\
  &= [\mathrm{Pa}]
\end{align*}

$\Rightarrow$ L'équation est homogène. ✓

\subsection*{Exercice 2 — Conversion de débits}

\textbf{Énoncé} : Une pompe de circulation a un débit nominal $Q = 2\,400~\mathrm{L/min}$.
Exprimer ce débit en m³/h et en m³/s.

\textbf{Solution} :

\begin{align*}
  Q &= 2\,400 \text{ L/min} \times \frac{1~\mathrm{m^3}}{1\,000~\mathrm{L}} \times
  \frac{60~\mathrm{min}}{1~\mathrm{h}} \\
  &= 2\,400 \times \frac{60}{1\,000}~\mathrm{m^3/h} \\
  &= 144~\mathrm{m^3/h}
\end{align*}

Puis :

\begin{align*}
  Q &= 144~\mathrm{m^3/h} \times \frac{1~\mathrm{h}}{3\,600~\mathrm{s}} \\
  &= 0{,}04~\mathrm{m^3/s}
\end{align*}

\subsection*{Exercice 3 — Résolution d'équation}

\textbf{Énoncé} : On mesure une consommation énergétique annuelle $E = 45~\mathrm{kWh}$. La
puissance moyenne est estimée à $\bar{P} = 3{,}5~\mathrm{kW}$. Combien de temps (en heures)
l'équipement a-t-il fonctionné ?

\textbf{Solution} :

Relation : $E = \bar{P} \cdot t$

Isoler $t$ :

\[
  t = \frac{E}{\bar{P}} = \frac{45}{3{,}5} = 12{,}86~\mathrm{h}
\]

Soit environ \textbf{12 h 52 min}.

\subsection*{Exercice 4 — Interpolation sur un abaque}

\textbf{Énoncé} : Soit l'abaque des propriétés de l'eau (pression atm.) :

\begin{center}
  \sffamily
  \begin{tabular}{|c|c|c|}
    \hline
    $T\,[\mathrm{°C}]$ & $\rho\,[\mathrm{kg/m^3}]$ & $c_p\,[\mathrm{kJ/(kg \cdot K)}]$ \\
    \hline
    20 & 998{,}2 & 4{,}181 \\
    40 & 992{,}2 & 4{,}179 \\
    60 & 983{,}2 & 4{,}185 \\
    \hline
  \end{tabular}
\end{center}

Déterminer $\rho$ et $c_p$ à $T = 50°C$ par interpolation linéaire.

\textbf{Solution} :

Pour $\rho$ (entre 40°C et 60°C) :

\begin{align*}
  \rho(50) &= 992{,}2 + (983{,}2 - 992{,}2) \times \frac{50 - 40}{60 - 40} \\
  &= 992{,}2 + (-9) \times \frac{10}{20} \\
  &= 992{,}2 - 4{,}5 \\
  &= 987{,}7~\mathrm{kg/m^3}
\end{align*}

Pour $c_p$ (entre 40°C et 60°C) :

\begin{align*}
  c_p(50) &= 4{,}179 + (4{,}185 - 4{,}179) \times \frac{50 - 40}{60 - 40} \\
  &= 4{,}179 + 0{,}006 \times 0{,}5 \\
  &= 4{,}182~\mathrm{kJ/(kg \cdot K)}
\end{align*}

% ============================================================
\section{Diagrammes et schémas de synthèse}
% ============================================================

\subsection{Les sept dimensions fondamentales du SI}

La figure~\ref{fig:dimensions-si} présente les sept grandeurs de base du Système International
et leurs unités associées.

\begin{figure}[h]
  \centering
  \begin{tikzpicture}[
      node distance = 1.4cm and 1.8cm,
      dim/.style  = {draw, rounded corners=4pt, minimum width=3.2cm, minimum height=0.9cm,
                     font=\sffamily\small, align=center, thick},
      base/.style = {dim, fill=FedBlue!15, draw=FedBlue},
      center/.style = {draw, circle, fill=FedBlue, text=white, font=\sffamily\bfseries\small,
                       minimum size=1.8cm, align=center},
      arr/.style  = {-Stealth, thick, FedBlue!70}
    ]
    % Centre
    \node[center] (SI) {SI\\Dimensions};

    % 7 nœuds autour
    \node[base, above=1.6cm of SI]               (M) {\textbf{Masse} $[M]$\\\si{\kilogram} (kg)};
    \node[base, above right=1.1cm and 2.2cm of SI](L) {\textbf{Longueur} $[L]$\\\si{\meter} (m)};
    \node[base, right=2.8cm of SI]               (T) {\textbf{Temps} $[T]$\\\si{\second} (s)};
    \node[base, below right=1.1cm and 2.2cm of SI](I) {\textbf{Courant} $[I]$\\\si{\ampere} (A)};
    \node[base, below=1.6cm of SI]               (th){\textbf{Température} $[\Theta]$\\\si{\kelvin} (K)};
    \node[base, below left=1.1cm and 2.2cm of SI] (N) {\textbf{Quantité de matière} $[N]$\\\si{\mole} (mol)};
    \node[base, left=2.8cm of SI]                (J) {\textbf{Intensité lumineuse} $[J]$\\\si{\candela} (cd)};

    % Flèches
    \foreach \n in {M,L,T,I,th,N,J}
      \draw[arr] (SI) -- (\n);
  \end{tikzpicture}
  \caption{Les sept dimensions fondamentales du Système International d'unités}
  \label{fig:dimensions-si}
\end{figure}

\subsection{Méthode d'homogénéité — Schéma de raisonnement}

La figure~\ref{fig:methode-homogeneite} illustre la démarche pas-à-pas pour vérifier
l'homogénéité d'une équation physique (exemple : bilan thermique $\dot{Q} = \dot{m}\,c_p\,\Delta T$).

\begin{figure}[h]
  \centering
  \begin{tikzpicture}[
      node distance=0.55cm,
      step/.style  = {draw, rounded corners=4pt, fill=FedBlue!12, draw=FedBlue,
                       text width=8cm, minimum height=0.85cm, align=left,
                       font=\sffamily\small, inner sep=6pt},
      result/.style= {draw, rounded corners=4pt, fill=FedGreen!20, draw=FedGreen,
                       text width=8cm, minimum height=0.85cm, align=left,
                       font=\sffamily\small, inner sep=6pt},
      err/.style   = {draw, rounded corners=4pt, fill=FedOrange!20, draw=FedOrange,
                       text width=8cm, minimum height=0.85cm, align=left,
                       font=\sffamily\small, inner sep=6pt},
      arr/.style   = {-Stealth, thick, FedBlue}
    ]
    \node[step]   (s1) {\textbf{Étape 1} — Écrire l'équation : $\dot{Q} = \dot{m}\cdot c_p\cdot\Delta T$};
    \node[step, below=of s1] (s2) {\textbf{Étape 2} — Remplacer chaque grandeur par sa dimension :
      $[\dot{Q}]=\mathrm{W}=\mathrm{kg\cdot m^2\cdot s^{-3}}$};
    \node[step, below=of s2] (s3) {\textbf{Étape 3} — Développer le membre droit :
      $[\dot{m}][c_p][\Delta T] = (\mathrm{kg\cdot s^{-1}})(\mathrm{m^2\cdot s^{-2}\cdot K^{-1}})(\mathrm{K})$};
    \node[step, below=of s3] (s4) {\textbf{Étape 4} — Simplifier et regrouper les exposants};
    \node[result, below=of s4] (ok) {\textcolor{FedGreen!60!black}{\textbf{Résultat :}}
      $\mathrm{kg\cdot m^2\cdot s^{-3} = kg\cdot m^2\cdot s^{-3}}$ \quad \textbf{Équation homogène} \checkmark};

    \draw[arr] (s1) -- (s2);
    \draw[arr] (s2) -- (s3);
    \draw[arr] (s3) -- (s4);
    \draw[arr] (s4) -- (ok);

    % Annotation à droite de s4
    \node[err, right=1.2cm of s4, text width=4.2cm]
      (warn) {\textcolor{FedOrange!80!black}{\textbf{Si différent :}} équation fausse — vérifier coefficients ou formule};
    \draw[-Stealth, dashed, FedOrange, thick] (s4.east) -- (warn.west);
  \end{tikzpicture}
  \caption{Démarche de vérification d'homogénéité dimensionnelle (exemple thermique)}
  \label{fig:methode-homogeneite}
\end{figure}

\subsection{Grandeurs physiques courantes en CVC — Rappel visuel}

\begin{figure}[h]
  \centering
  \begin{tikzpicture}[
      col/.style={draw=FedBlue, fill=FedBlue!8, rounded corners=3pt,
                  minimum width=2.8cm, minimum height=0.75cm,
                  font=\sffamily\footnotesize, align=center},
      hdr/.style={draw=FedBlue, fill=FedBlue, text=white, rounded corners=3pt,
                  minimum width=2.8cm, minimum height=0.8cm,
                  font=\sffamily\small\bfseries, align=center},
      colO/.style={col, draw=FedOrange, fill=FedOrange!10},
      hdrO/.style={hdr, draw=FedOrange, fill=FedOrange},
      colG/.style={col, draw=FedGreen, fill=FedGreen!10},
      hdrG/.style={hdr, draw=FedGreen, fill=FedGreen},
      node distance=0pt
    ]
    % Colonne 1 — Thermique
    \node[hdr]  (h1) at (0,0)        {Thermique};
    \node[col, below=of h1] (r11)    {Puissance : \si{\watt} (W)};
    \node[col, below=of r11] (r12)   {Énergie : \si{\kilo\watt\hour} (kWh)};
    \node[col, below=of r12] (r13)   {Temp. : \si{\kelvin} / \si{\degreeCelsius}};
    \node[col, below=of r13] (r14)   {$c_p$ : J/(kg·K)};

    % Colonne 2 — Hydraulique
    \node[hdrO] (h2) at (3.2,0)      {Hydraulique};
    \node[colO, below=of h2] (r21)   {Pression : \si{\pascal} (Pa)};
    \node[colO, below=of r21] (r22)  {Débit vol. : m³/h};
    \node[colO, below=of r22] (r23)  {Débit mass. : kg/s};
    \node[colO, below=of r23] (r24)  {Vitesse : m/s};

    % Colonne 3 — Aéraulique / géométrie
    \node[hdrG] (h3) at (6.4,0)      {Géométrie/Air};
    \node[colG, below=of h3] (r31)   {Longueur : \si{\meter} (m)};
    \node[colG, below=of r31] (r32)  {Surface : \si{\meter\squared} (m²)};
    \node[colG, below=of r32] (r33)  {Volume : \si{\meter\cubed} (m³)};
    \node[colG, below=of r33] (r34)  {Masse vol. : kg/m³};
  \end{tikzpicture}
  \caption{Grandeurs physiques et unités courantes en génie climatique (CVC)}
  \label{fig:grandeurs-cvc}
\end{figure}

% ============================================================
\section{Données terrain — Extraits de catalogues fabricants}
% ============================================================

Ce chapitre prépare l'étudiant à exploiter des données techniques réelles. Les extraits
ci-dessous sont représentatifs des catalogues industriels Grundfos (circulateurs) et
Daikin (pompes à chaleur) — compétence \textbf{C1-1 : exploiter des sources d'information}.

\subsection{Extrait catalogue Grundfos — Circulateur MAGNA3 25-100}

La gamme \textbf{Grundfos MAGNA3} est représentative des circulateurs haute efficacité
(classe IE5) utilisés en chauffage et climatisation. Les données suivantes correspondent
aux conditions d'essai EN\,1151 (eau à \SI{60}{\celsius}, $\rho = \SI{983}{\kilogram\per\metre\cubed}$).

\begin{table}[H]
  \centering\sffamily\small
  \caption{Extrait performances Grundfos MAGNA3 25-100 — Points caractéristiques (eau 60\,°C, EN\,1151)}
  \label{tab:grundfos-magna3}
  \begin{tabular}{|c|c|c|c|c|}
    \hline
    \textbf{Débit} $\dot{V}$ & \textbf{HMT} & \textbf{Puiss. absorbée} $P_1$ & \textbf{Rendement} $\eta$ & \textbf{Vitesse} $n$ \\
    (\si{\metre\cubed\per\hour}) & (\si{\meter} CE) & (\si{\watt}) & (\%) & (tr/min) \\
    \hline
    0{,}0 & 10{,}0 &  85 &  0{,}0 & 3\,200 \\
    0{,}5 &  9{,}5 &  92 & 14{,}3 & 3\,100 \\
    1{,}0 &  8{,}6 & 105 & 22{,}6 & 2\,950 \\
    1{,}5 &  7{,}4 & 118 & 25{,}8 & 2\,750 \\
    2{,}0 &  6{,}0 & 128 & 25{,}8 & 2\,550 \\
    2{,}5 &  4{,}5 & 132 & 23{,}5 & 2\,350 \\
    3{,}0 &  2{,}6 & 128 & 16{,}9 & 2\,100 \\
    3{,}2 &  1{,}8 & 120 & 13{,}2 & 2\,000 \\
    \hline
  \end{tabular}
\end{table}

\textbf{Rappel — Conversion HMT $\leftrightarrow$ pression différentielle :}
\begin{equation}
  \Delta P\,[\si{\pascal}] = \rho \cdot g \cdot H_{MT}\,[\si{\meter}]
  \quad\Rightarrow\quad
  \Delta P \approx 983 \times 9{,}81 \times H_{MT} \approx 9\,644 \times H_{MT}
  \label{eq:hmt-dp}
\end{equation}

\begin{exemple}[title={Application — Lecture catalogue Grundfos MAGNA3}]
\textbf{Contexte :} Un circuit de chauffage présente une perte de charge totale
$\Delta P_{circuit} = \SI{45\,000}{\pascal}$ pour un débit nominal $\dot{V} = \SI{1{,}5}{\metre\cubed\per\hour}$.

\textbf{1 — Conversion en mètres de colonne d'eau (vérification dimensionnelle) :}
\[
  H_{circuit} = \frac{\Delta P}{\rho \cdot g}
  = \frac{45\,000\,[\si{\pascal}]}{983\,[\si{\kilogram\per\metre\cubed}] \times 9{,}81\,[\si{\metre\per\second\squared}]}
  = \frac{45\,000}{9\,644} \approx \SI{4{,}67}{\meter}~\text{CE}
\]

\textbf{Homogénéité :} $\dfrac{\mathrm{Pa}}{\mathrm{kg/m^3 \cdot m/s^2}}
= \dfrac{\mathrm{kg \cdot m^{-1} \cdot s^{-2}}}{\mathrm{kg \cdot m^{-2} \cdot s^{-2}}}
= \mathrm{m}$ \checkmark

\textbf{2 — Vérification du point de fonctionnement :}

À $\dot{V} = 1{,}5\,\mathrm{m^3/h}$, la MAGNA3 25-100 développe $H_{MT} = 7{,}4\,\mathrm{m}$ CE.
La pompe peut alimenter ce circuit ($7{,}4 > 4{,}67\,\mathrm{m}$).
Le point de fonctionnement réel (intersection courbe pompe / résistance réseau) se situe
entre $\dot{V} = 1{,}5$ et $\SI{2{,}0}{\metre\cubed\per\hour}$.

\textbf{3 — Rendement au point de fonctionnement :}
$\eta \approx 25{,}8\,\%$ — valeur typique d'un petit circulateur de chauffage.
\end{exemple}

\subsection{Extrait catalogue Daikin — Pompe à chaleur Altherma 3 R (air/eau)}

La gamme \textbf{Daikin Altherma 3 R} est représentative des PAC air/eau résidentielles.
Les COP sont mesurés selon EN\,14511 ; ils \textbf{diminuent} quand la température extérieure
baisse ou que la température de départ augmente.

\begin{table}[H]
  \centering\sffamily\small
  \caption{Extrait performances Daikin Altherma 3 R 8\,kW — COP selon conditions d'essai EN\,14511}
  \label{tab:daikin-cop}
  \begin{tabular}{|l|c|c|c|c|}
    \hline
    \textbf{Conditions} & $T_{ext}$ & $T_{départ}$ & \textbf{Puiss. calorifique} & \textbf{COP} \\
    & (\si{\degreeCelsius}) & (\si{\degreeCelsius}) & (\si{\kilo\watt}) & (—) \\
    \hline
    A7/W35 (nominale)      & $+7$  & 35 & 8{,}0 & 4{,}65 \\
    A7/W45                 & $+7$  & 45 & 8{,}1 & 3{,}40 \\
    A2/W35 (intermédiaire) & $+2$  & 35 & 7{,}3 & 3{,}80 \\
    A$-$7/W35 (hiver froid)  & $-7$  & 35 & 6{,}0 & 2{,}90 \\
    A$-$10/W35 (hiver froid) & $-10$ & 35 & 5{,}4 & 2{,}60 \\
    A$-$15/W35 (grand froid) & $-15$ & 35 & 4{,}5 & 2{,}20 \\
    A7/W55 (haute temp.)   & $+7$  & 55 & 8{,}5 & 2{,}55 \\
    \hline
  \end{tabular}
\end{table}

\textbf{Clé de lecture :} la notation \textbf{A7/W35} signifie
$T_{air\_ext} = +7\,\si{\degreeCelsius}$ / $T_{eau\_départ} = 35\,\si{\degreeCelsius}$.

\begin{exemple}[title={Application — Vérification dimensionnelle avec données Daikin}]
\textbf{Contexte :} Maison de $120\,\mathrm{m^2}$ à Strasbourg (zone H1c,
$T_{ext,base} = \SI{-11}{\celsius}$, $DJU = 3\,100$).
Besoin de pointe : $\Phi_{ch} = \SI{6{,}0}{\kilo\watt}$.

\textbf{1 — Puissance électrique absorbée en pointe ($T_{ext} = -15\,\si{\celsius}$, COP = 2,20) :}
\[
  P_{el,pointe} = \frac{\Phi_{ch}}{COP}
  = \frac{6{,}0\,[\si{\kilo\watt}]}{2{,}20\,[-]}
  = \SI{2{,}73}{\kilo\watt}
  \quad\text{Homogénéité :}\ \frac{\mathrm{kW}}{1} = \mathrm{kW}\ \checkmark
\]

À $-15\,\si{\celsius}$, l'Altherma 3R 8\,kW ne délivre que $4{,}5\,\mathrm{kW}$
$< 6{,}0\,\mathrm{kW}$ requis. \textbf{Appoint électrique nécessaire} ou modèle 12\,kW.

\textbf{2 — Consommation annuelle ($SCOP_{annuel} = 3{,}6$) :}
\[
  W_{el,annuel} = \frac{\Phi_{ch} \times DJU \times 24}{SCOP \times \Delta T_{base}}
  = \frac{6{,}0 \times 3\,100 \times 24}{3{,}6 \times 30}
  = \frac{446\,400}{108}
  \approx \SI{4\,133}{\kilo\watt\hour\per\year}
\]

\textbf{Homogénéité :} $\dfrac{[\si{\kilo\watt}] \times [\text{jours}] \times [\si{\hour\per\day}]}{[1] \times [\si{\kelvin}]}
= \dfrac{\mathrm{kW \cdot h}}{\mathrm{K} / \mathrm{K}} = \mathrm{kWh/an}$ \checkmark
\end{exemple}

\subsection*{Application — Schéma d'installation hydraulique avec symboles ISO GCF}

La figure~\ref{fig:installation-dim-gcf} illustre un circuit de chauffage bitubes
représenté avec les \textbf{symboles normalisés ISO GCF} des composants et annoté
avec les grandeurs physiques mobilisées lors de l'analyse dimensionnelle.

\tikzset{
  pics/pompe/.code={
    \draw[FedOrange, thick] (0,0) circle (0.42);
    \fill[FedOrange!55] (0,0)--(90:0.37)  arc(90:30:0.37) --cycle;
    \fill[FedOrange!55] (0,0)--(210:0.37) arc(210:150:0.37)--cycle;
    \fill[FedOrange!55] (0,0)--(330:0.37) arc(330:270:0.37)--cycle;
    \fill[FedOrange]    (0,0) circle (0.07);
  },
  pics/chaudiere/.code={
    \draw[FedRed, thick] (-0.55,-0.40) rectangle (0.55,0.40);
    \draw[FedRed, line width=1.4pt]
      (-0.10,-0.35)..controls(-0.25,0.05)and(-0.05,0.20)..(0,0.35)
      ..controls(0.05,0.20)and(0.25,0.05)..(0.10,-0.35);
  },
  pics/vase-expansion/.code={
    \draw[FedGreen, thick] (0,0) circle (0.38);
    \draw[FedGreen, thick] (-0.38,0)--(0.38,0);
    \fill[FedGreen!25] (-0.38,0) arc(180:360:0.38)--cycle;
  },
  pics/radiateur/.code={
    \draw[FedRed!75, thick] (-0.45,-0.35) rectangle (0.45,0.35);
    \foreach \x in {-0.22,0.0,0.22}{
      \draw[FedRed!50, thin] (\x,-0.35)--(\x,0.35);
    }
  },
}

\begin{figure}[H]
  \centering
  \begin{tikzpicture}[thick, >=Stealth, font=\small]

    % ---- Équipements — symboles ISO GCF ----
    \pic at (0, 0)     {chaudiere};
    \pic at (0,-2.8)   {pompe};
    \pic at (5.5, 0)   {radiateur};
    \pic at (2.75,1.2) {vase-expansion};

    % ---- Libellés équipements ----
    \node[font=\scriptsize\bfseries, text=FedRed]
      at (0,-0.65)    {Chaudière};
    \node[font=\scriptsize\bfseries, text=FedOrange]
      at (-0.8,-2.8)  {Pompe};
    \node[font=\scriptsize\bfseries, text=FedRed!80]
      at (5.5,-0.65)  {Radiateur};
    \node[font=\scriptsize\bfseries, text=FedGreen!70!black]
      at (2.75,1.80)  {Vase exp.};

    % ---- Tuyauterie départ (rouge) ----
    \draw[->, FedRed, very thick]
      (0.55,0) -- node[above, font=\scriptsize, text=FedRed]
        {$T_{dep}=70^\circ$C \quad $\dot{V}$\,[m³/h]}
      (5.05,0);

    % ---- Raccordement vase d'expansion sur le départ ----
    \draw[FedGreen, thick] (2.75,0)--(2.75,0.82);

    % ---- Tuyauterie retour (bleue) : radiateur → collecteur → pompe ----
    \draw[FedBlue, very thick] (5.5,-0.35)--(5.5,-4.2)--(0,-4.2)--(0,-3.22);

    % ---- Pompe → chaudière (sortie pompe) ----
    \draw[->, FedBlue, very thick]
      (0,-2.38) -- node[right, font=\scriptsize, text=FedBlue]
        {$T_{ret}=55^\circ$C}
      (0,-0.40);

    % ---- Encart formule bilan et grandeurs physiques ----
    \node[draw=FedBlue!35, fill=FedBlue!5, rounded corners=3pt,
          font=\scriptsize, inner sep=6pt, align=left, text width=6.5cm]
      at (3.2,-5.4) {
        \textbf{Bilan thermique :} $\dot{Q}=\dot{V}\cdot\rho\cdot c_p\cdot\Delta T$
        \hfill$[\si{\watt}]$\\[3pt]
        \textbf{Grandeurs :}
        $\dot{V}$\,[m³/h] \quad $\Delta P$\,[\si{\pascal}] \quad
        $\rho$\,[kg/m³] \quad $\dot{m}$\,[kg/s]
      };

  \end{tikzpicture}
  \caption{Circuit de chauffage bitubes — symboles ISO GCF et grandeurs physiques pour
           l'analyse dimensionnelle}
  \label{fig:installation-dim-gcf}
\end{figure}

% ============================================================
\section{Synthèse et bonnes pratiques}
% ============================================================

\begin{tcolorbox}[colback=FedBlue!10, colframe=FedBlue, coltitle=white, title=Checklist — Analyse dimensionnelle et conversions]
  Avant d'appliquer une formule ou un résultat de calcul :

  \begin{itemize}
    \item ✓ Vérifier l'homogénéité dimensionnelle de l'équation
    \item ✓ S'assurer que toutes les variables sont exprimées dans les bonnes unités
    \item ✓ Convertir les données d'entrée si nécessaire
    \item ✓ Effectuer le calcul
    \item ✓ Exprimer le résultat dans l'unité requise
    \item ✓ Valider l'ordre de grandeur du résultat
    \item ✓ Justifier chaque étape de la transformation d'équation pour éviter les erreurs
  \end{itemize}
\end{tcolorbox}

\finchapitre
